Ce projet étudie la gravité des blessures des usagers impliqués dans les accidents
corporels de la route enregistrés en France métropolitaine, Corse comprise. Il poursuit
deux objectifs complémentaires~: décrire les différences de gravité observées selon les
caractéristiques des usagers et le contexte de l'accident, puis évaluer la capacité d'une
régression logistique à estimer cette gravité sur une année ultérieure.

L'étude s'appuie sur les données du Bulletin d'analyse des accidents corporels de la
circulation, ou BAAC, pour les années 2019 à 2024. Les quatre fichiers annuels décrivent
les caractéristiques générales de l'accident, les lieux, les véhicules et les usagers.
L'unité d'analyse retenue est la participation d'un usager à un accident. La cible
principale distingue les usagers décédés ou hospitalisés des usagers indemnes ou
légèrement blessés~; les gravités inconnues restent identifiées séparément. Ce
regroupement constitue une convention d'analyse, avec les limites propres à la collecte
de l'hospitalisation. BAAC ne décrivant pas l'ensemble des déplacements, le projet
n'estime pas un risque d'accident par trajet.

La première partie porte sur la compréhension, la préparation et l'analyse des données.
L'audit des six années met en évidence des différences de schéma, des codes particuliers,
des informations manquantes et plusieurs enregistrements de lieux pour certains accidents.
Chaque difficulté fait l'objet d'une règle justifiée et de contrôles associés. Les
jointures préservent les usagers, les valeurs sources restent traçables et les désaccords
entre lieux sont signalés. La table préparée comprend ainsi 702\,823 observations
d'usagers, correspondant à 309\,756 accidents métropolitains.

Les analyses exploratoires et statistiques se concentrent sur 2019-2022, la période
d'apprentissage. Parmi les 467\,079 usagers dont la gravité est connue sur cette période,
environ 17,8\,\% sont décédés ou hospitalisés. Les comparaisons font notamment apparaître
des proportions plus élevées chez les piétons que chez les conducteurs, chez les personnes
de 75 ans ou plus que chez les 35-44 ans, de nuit sans éclairage public qu'en plein jour,
et hors agglomération qu'en agglomération. Les écarts et rapports de proportions sont
accompagnés d'intervalles bootstrap obtenus en rééchantillonnant des accidents entiers,
séparément dans chaque année. Une lecture annuelle et une analyse des décès seuls
complètent ces résultats. Ces associations restent exploratoires et non causales.

La seconde partie compare une référence constante à une régression logistique L2,
avec onze variables et un encodage appris sur 2019-2022. La validation sur 2023 retient
$C=1$ ; cette configuration est évaluée sur 2024 sans ajout de 2023 à l'apprentissage.
Sur les 117\,278 usagers du test, la log-loss passe de 0,4703 à 0,4001, soit une réduction
relative de 14,93\,\%. La ROC-AUC atteint 0,770. Le gain se maintient donc sur l'année
suivante, mais le rappel au seuil illustratif de 0,5 reste limité à 18,22\,\%.

Les probabilités élevées sont surestimées et le classement est moins performant chez
les piétons. Le signal porté par les informations manquantes impose aussi une lecture
prudente. Cette étude rétrospective met en évidence un apport prédictif mesurable,
sans démontrer une utilité opérationnelle ni un effet causal. Les données, les règles,
les résultats et leurs limites sont documentés dans les notebooks et le rapport.

\vspace{3mm}
{\sffamily\small\raggedright\textbf{Mots-clés~:} BAAC, accidents corporels, gravité des blessures,
analyse de données, bootstrap, apprentissage automatique.\par}
